\documentclass[10pt,a4paper]{amsart}
\usepackage[margin=19mm]{geometry}
\usepackage{amsmath,amssymb,mathtools}
\usepackage[scr=rsfs,cal=euler]{mathalfa}
\usepackage{xfrac}
\usepackage[unicode,hidelinks,bookmarks=false]{hyperref}
\newcommand{\ie}{i.e.\ }
\newcommand{\cf}{cf.\ }
\newcommand{\resp}{respectively\ }
\newcommand{\Subsection}{\subsection}
\providecommand{\nobreakdash}{\nobreak\hbox{-}}
\setlength{\emergencystretch}{2em}
\multlinegap0pt
\usepackage{bm}
\usepackage{bbm}
\usepackage{dsfont}
\usepackage{xspace}
\usepackage{cancel}
\usepackage{mathtools}
\usepackage{amsthm}
\makeatletter
\@ifundefined{theorem}{\newtheorem{theorem}{Theorem}}{}
\@ifundefined{lemma}{\newtheorem{lemma}[theorem]{Lemma}}{}
\makeatother
\makeatletter
\expandafter\ifx\csname claim\endcsname\relax
\newenvironment{claim}[1]{\par\noindent\underline{\emph{Claim}:}\space#1}{}
\fi
\expandafter\ifx\csname claimproof\endcsname\relax
\newenvironment{claimproof}[1]{\par\noindent{\emph{Proof of Claim}:}\space#1}{\hfill $\blacksquare$}
\fi
\newcommand{\cof}{\mathrm{cf}}
\newcommand{\dom}{\mathrm{dom}}
\newcommand{\mbb}{\mathbb}
\newcommand{\mcal}{\mathcal}
\makeatother
\title[On $2$-stationarity of $\mathcal{P}\_{\kappa}\lambda$]{On $2$-stationarity of $\mathcal{P}\_{\kappa}\lambda$}
\author{Hiroshi Sakai, M. Catalina Torres}
\date{Source: arXiv:2607.16563v1 (CC BY 4.0)}
\begin{document}
\maketitle

\noindent\emph{Source and licence.} On $2$-stationarity of $\mathcal{P}_{\kappa}\lambda$. Version arXiv:2607.16563v1 (CC BY 4.0), distributed under the Creative Commons Attribution 4.0 International licence, as recorded in the arXiv licence metadata for this exact version and shown on the abstract page https://arxiv.org/abs/2607.16563v1. Full rights notice: licenses/arxiv-2607.16563-CC-BY-4.0.txt.\par\smallskip

\noindent\textit{Editorial scope.} Counted unit: the Lemma whose source label is \texttt{lem:Q\_basics} in the article (source0.tex lines 439--445, source label \texttt{lem:Q\_basics}), delivered together with the article's own complete original proof of it (source0.tex lines 447--484, one \texttt{proof} environment of 38 lines). No mathematics is written, repaired or re-derived: statement and proof are the article's own text line for line. Required context retained uncounted: none. Every definition, display and label of those lines is retained unchanged. Omitted from this package and not redistributed: 1--438, 446--446, 485--1084. Nothing inside a retained excerpt is omitted or altered: every retained body is delivered exactly as the article's lines are written, so no omission receipt of any kind accompanies this package. Citations: the retained excerpts carry no citation command at all, so no bibliography entry of the article is transcribed and no reference is redistributed. Reference adaptation: none was needed and none was made; every reference inside the delivered unit resolves inside it, so no unresolved reference or citation is produced. Printed slips kept literal: no printed word, symbol, hypothesis or number of the counted statement or of its proof was altered. Layout and class: the article's own class is \texttt{amsart} with packages including babel, inputenc, fontenc, amssymb, mathtools, bbm, geometry, tcolorbox; the wrapper is the lane's pinned \texttt{amsart} editorial wrapper at 10pt on A4, which restates the theorem-like environments the retained text needs and adds the article's own macro definitions that the retained text needs (\texttt{\textbackslash cof}, \texttt{\textbackslash dom}, \texttt{\textbackslash mbb}, \texttt{\textbackslash mcal}), with extra wrapper packages (bm, bbm, dsfont, xspace, cancel, mathtools). The delivered numbering is the unit's own. Assets: no retained span contains a figure environment, a graphics-inclusion command, a PDF-page inclusion, a table or a TikZ picture; no image file is copied into this package, the package declares no asset dependency and no image file is redistributed with it. Licence: version arXiv:2607.16563v1 of this article is distributed under the Creative Commons Attribution 4.0 International licence, as recorded in the arXiv licence metadata for this exact version and shown on the abstract page https://arxiv.org/abs/2607.16563v1. A full rights notice is delivered at licenses/arxiv-2607.16563-CC-BY-4.0.txt. Characters: every retained excerpt is pure ASCII, so the delivered engine, XeLaTeX with its default Latin Modern text font, reports no missing character.
\begin{lemma} \label{lem:Q_basics}
Let $\kappa$ and $\nu$ be regular uncountable cardinals with $\kappa \leq \nu$.
\begin{enumerate}
\item For any $\gamma < \nu$ the set $D_\gamma = \{ p \in \mbb{Q} ( \kappa , \nu ) : \gamma \leq \delta_p \}$ is dense in $\mbb{Q} ( \kappa , \nu )$.
\item $\mbb{Q} ( \kappa , \nu )$ is $\nu$-strategically closed.
\end{enumerate}
\end{lemma}

\begin{proof}
Let $\mbb{Q} := \mbb{Q} ( \kappa , \nu )$.

\medskip
\noindent
(1) Suppose $\gamma < \nu$ and $p \in \mbb{Q}$. Take $\delta$ with $\gamma , \delta_p \leq \delta < \nu$, and let $q : \mcal{P}_\kappa \delta \to 2$ be such that $q \supseteq p$, and $q(z) = 0$ for all $z \in \mcal{P}_\kappa \delta \setminus \dom (p)$. It is easy to check that $q \in \mbb{Q}$. Then $q \leq p$ and $q \in D_\gamma$.

\medskip
\noindent
(2) First we make some preliminaries.

Suppose $\alpha$ is a limit ordinal $< \nu$, and $\vec{p} = \langle p_\beta : \beta < \alpha \rangle$ is a descending sequence in $\mbb{Q}$. Then let $\delta_{\vec{p}} := \sup_{\beta < \alpha} \delta_{p_\beta}$, $q_{\vec{p}} := \bigcup_{\beta < \alpha} p_\beta$ and $r_{\vec{p}} : \mcal{P}_\kappa \delta_{\vec{p}} \to 2$ be such that $r_{\vec{p}} \supseteq q_{\vec{p}}$ and $r_{\vec{p}} (z) = 0$ for all $z \in \mcal{P}_\kappa \delta_{\vec{p}} \setminus \dom ( q_{\vec{p}} )$.

Note that if $\vec{p}$ is not eventually constant (i.e.~$\delta_{p_\beta} < \delta_{\vec{p}}$ for all $\beta < \alpha$), and $\cof ( \delta_{\vec{p}} ) < \kappa$, then $\mcal{P}_\kappa \delta_{\vec{p}} \setminus \dom ( q_{\vec{p}} )$ is the set of all $z \in \mcal{P}_\kappa \delta_{\vec{p}}$ with $\sup z = \delta_{\vec{p}}$. Otherwise, $\mcal{P}_\kappa \delta_{\vec{p}} \setminus \dom ( q_{\vec{p}} ) = \emptyset$, and so $r_{\vec{p}} = q_{\vec{p}}$.
Note also that if $\vec{p}$ has a lower bound in $\mbb{Q}$, then $r_{\vec{p}} \in \mbb{Q}$. This is because ${r_{\vec{p}}}^{-1} (1) \subseteq p^{-1} (1)$ for any lower bound $p$ of $\vec{p}$.

Now we give a strategy of Even in the game $\Game_\nu ( \mbb{Q} )$. Suppose $\alpha$ is an even ordinal $< \nu$, and $\vec{p} = \langle p_\beta : \beta < \alpha \rangle$ has been played before the $\alpha$-th stage. If $\alpha$ is a successor ordinal, then Even chooses $p_\alpha$ so that $p_\alpha \lneq p_{\alpha -1}$. If $\alpha$ is a limit ordinal, and $\vec{p}$ has a lower bound in $\mbb{Q}$, then Even chooses $p_\alpha = r_{\vec{p}}$.

We show that the above is a winning strategy of Even. Suppose $\alpha$ is a limit ordinal $< \nu$, and $\vec{p} = \langle p_\beta : \beta < \alpha \rangle$ is a play of $\Game_\nu ( \mbb{Q} )$ in which Even has moved according to the above strategy. We must show that $\vec{p}$ has a lower bound in $\mbb{Q}$. For this it suffices to prove that $r_{\vec{p}} \in \mbb{Q}$.

Before starting, note the following by the strategy of Even at successor stages.
\begin{itemize}
\item[(i)] $p_\gamma \lneq p_{\gamma -1}$, and so $\delta_{p_\gamma} > \delta_{p_{\gamma -1}}$, for all even successor $\gamma < \alpha$.
\end{itemize}
Below, let $\delta$, $q$ and $r$ denote $\delta_{\vec{p}}$, $q_{\vec{p}}$ and $r_{\vec{p}}$, respectively.

Suppose $x \in \mcal{P}_\kappa \nu$, and $\mu := x \cap \kappa$ is a regular uncountable cardinal. We show that $r^{-1} (1) \cap \mcal{P}_\mu x$ is non-stationary. This is clear if $x \not\subseteq \delta$. So assume $x \subseteq \delta$. We consider several cases. Note that $\delta$ is a limit ordinal by (i).

First suppose $\sup x < \delta$. Then there is $\beta < \alpha$ with $x \in \mcal{P}_\kappa \delta_{p_\beta}$. Then $r^{-1} (1) \cap \mcal{P}_\mu x = {p_\beta}^{-1} (1) \cap \mcal{P}_\mu x$, and it is non-stationary since $p_\beta \in \mbb{Q}$.

Next suppose $\sup x = \delta$, and $\cof ( \delta ) < \mu$. Then there are club many $z \in \mcal{P}_\mu x$ with $\sup z = \delta$, and $r(z) = 0$ for all such $z$. So $r^{-1} (1) \cap \mcal{P}_\mu x$ is non-stationary.

Finally suppose $\sup x = \delta$, and $\cof ( \delta ) \geq \mu$. Let $C$ be the collection of $\delta_{p_\beta}$ for all limit $\beta < \alpha$. Then $C$ is club in $\delta$. So the set
\[
Z := \{ z \in \mcal{P}_\mu x : \sup z \notin z \wedge \sup z \in C \}
\]
is club in $\mcal{P}_\mu x$. But $r(z) = 0$ for all $z \in Z$, by (i) and the strategy of Even at limit stages. So $r^{-1} (1) \cap \mcal{P}_\mu x$ is non-stationary.
\end{proof}

\end{document}
